\clearpage
\chapter*{Conclusion générale}
\phantomsection
\addcontentsline{toc}{chapter}{Conclusion générale}
\markboth{Conclusion générale}{Conclusion générale}
\vspace{-1\baselineskip}
{\setstretch{1.40}

L'objectif de ce projet était de transformer un historique de ventes hétérogène en
un outil local capable d'aider concrètement les équipes de PAUL Maroc à décider
quoi produire et quoi commander. Pour atteindre cet objectif, nous avons construit
une chaîne complète qui consolide les ventes journalières et mensuelles, compare
dix candidats de prévision selon le produit, réconcilie les résultats avec un
niveau global, intègre les fêtes marocaines, les matchs, les événements et les
commandes clients, puis convertit les quantités recommandées en besoins de matières
premières grâce aux recettes et aux produits semi-finis. Le résultat prend la
forme de plans quotidien et mensuel sécurisés, d'un bon de commande,
d'indicateurs de coût et d'une interface accessible. La version étudiée
traite 373\,875 lignes de ventes, couvre plus de mille produits, sélectionne un
modèle pour 848 références et calcule 236 ingrédients distincts. La validation
mensuelle glissante obtient un WAPE de 8,25\,\% sur les quantités agrégées avec
Holt--Winters, tandis que les 86 tests automatisés réussissent sur le dernier
commit synchronisé. Ces résultats montrent que l'architecture répond aux besoins
fondamentaux de reproductibilité, de traçabilité et de confidentialité. Ils doivent
cependant être interprétés avec esprit critique. La précision journalière se
dégrade sur la longue traîne et la fenêtre récente révèle une sous-prévision
globale ; les ventes disponibles lors de l'audit s'arrêtent au 28 juin 2026 ; une
partie des recettes et des prix reste estimée ; enfin, le food-cost calculé ne
comprend pas encore toutes les charges. Le dashboard rend ces limites visibles
afin que l'utilisateur ne confonde pas une recommandation avec une certitude. Les
prochaines améliorations devraient donc porter en priorité sur l'alimentation SQL
continue, la validation systématique des recettes par le chef, l'intégration des
stocks réels et des délais fournisseurs, le suivi de la casse et des invendus, et
la mise en place d'alertes lorsque le biais dépasse un seuil. Une évolution vers
un stockage relationnel des configurations et des historiques d'exécution
faciliterait également le multi-site et la gestion des droits. Enfin, un modèle
local de langage pourrait enrichir l'assistant en langage naturel, à condition de
conserver les outils déterministes comme seule source de chiffres et de ne jamais
faire sortir les données de l'entreprise. Ce projet nous a permis d'articuler
ingénierie logicielle, données, séries temporelles et contraintes opérationnelles.
\par}
